% Nouveau EGA V (cote 25): archive cote 25, site page 127 (PDF page 128).
% English translation of the checked French annotated master.
% AI-translated by GPT-6 Astra (OpenAI, ChatGPT Pro mode).
\scanpage{25/127}
\mspage{\hand{30}}
braic closure of k(s), there exists a prime cycle T associated with $X_k$ such that
$k \hand{\to} H^0(T,\underline{O}_T)$ is an isomorphism (a condition automatically satisfied
when g is proper and surject\uline{i}ve). Then the morphism D: $\overset{\hand{\scriptscriptstyle\vee}}{P} \hand{\longrightarrow} \mathrm{Div}_{X/S}$ is
a monomorphism.

b) \var{\xout{Supposexxf}}{}{typewriter strike-out}{ega25_25_127-v0001} Consider the map $u \hand{\rightsquigarrow} D\circ u$ from
$\mathrm{Aut}_S(\overset{\hand{\scriptscriptstyle\vee}}{P}\ )$ to the set of families of divisors on X/S parametrized
by P . Then when \var{\over{f}{g}}{g}{writing over}{ega25_25_127-v0002} is surjective, the preceding map is injective,
in particular D = D u implies u = $\mathrm{id}_{\overset{\hand{\scriptscriptstyle\vee}}{P}}$\var{\inserted{; more generally, the morphism of functors $\underline{\mathrm{Aut}}_S(\overset{\hand{\scriptscriptstyle\vee}}{P}) \hand{\longrightarrow} \underline{\mathrm{Syst.\ lin.\ div}}_{X/S,\overset{\hand{\scriptscriptstyle\vee}}{P}}$ is a monomorphism}}{; more generally, the morphism of functors $\underline{\mathrm{Aut}}_S(\overset{\hand{\scriptscriptstyle\vee}}{P}) \hand{\longrightarrow} \underline{\mathrm{Syst.\ lin.\ div}}_{X/S,\overset{\hand{\scriptscriptstyle\vee}}{P}}$ is a monomorphism}{insertion}{ega25_25_127-v0003}

Observe that under the hypotheses of a), b) is a trivial consequence
of a); on the other hand, b) holds under less restrictive
hypotheses than a). Care must be taken that a) becomes false when one abandons
the restrictive hypothesis that was made there: take for example
S=Spec(k), X an open subset of $P^1_k$ not containing the two distinct points
a and b of $P^1_k(k)$ , D the family of divisors defined by the
canonical immersion of X into $P^1_k$ . Then the two \var{\xout{morphisms}}{}{typewriter strike-out}{ega25_25_127-v0005} points a and b
define the same divisors of X (namely the zero divisor), without being
identical themselves.

First suppose that S is the spectrum of a field \var{\inserted{k, which one can obviously (by ``descent'') assume to be}}{k, which one can obviously (by ``descent'') assume to be}{insertion}{ega25_25_127-v0006} alg closed. \var{\deleted{\illegible{1}}}{}{deletion}{ega25_25_127-v0007}
Let T be as in a)\var{\inserted{, endowed with the induced reduced structure,}}{, endowed with the induced reduced structure,}{insertion}{ega25_25_127-v0008} there is then an ``induced divisor'' morphism
$\mathrm{Div}_{X/k} \hand{\longrightarrow} \mathrm{Div}_{T/k}$ , and it suffices to prove that the composite $\overset{\hand{\scriptscriptstyle\vee}}{P} \hand{\longrightarrow} \mathrm{Div}_{T/k}$
is a monomorphism. Since the latter is again a lin syst of div,
one is reduced to the case X=T , hence to the case $H^0(X,\underline{O}_X) \hand{\xleftarrow{\sim}} k$ . Then for
every S over k, one has $g_{S*}(\underline{O}_{X_S}) \hand{\xleftarrow{\sim}} \underline{O}_S$ , hence (when $\underline{L}$ on X and u: $\underline{E} \hand{\to}$ \var{\deleted{\illegible{1}}}{}{deletion}{ega25_25_127-v0009}$\hand{g_*}(\underline{L})$
are as in \hand{4}.1.) when two everywhere nonzero sections $\hand{\phi}$, $\hand{\psi}$ of $\underline{E}_S$
are such that $u(\hand{\phi})$ and $u(\hand{\psi})$ are sections of $\underline{L}_S$ on $X_S$ having the same
divisor, hence \var{\xout{xxxxxf}}{}{typewriter strike-out}{ega25_25_127-v0010} obtained from each other by multiplication by
an \var{\inserted[typed]{invertible}}{invertible}{insertion}{ega25_25_127-v0011} section of $\underline{O}_{X_S}$, it follows that $\hand{\psi}$ is obtained from $\hand{\phi}$ by \unsure{multiplication}

\mspage{\hand{127}}
\var{\margin{Not automatically satisfied when the $X_s$ are \var{\deleted{\illegible{1}}}{}{deletion}{ega25_25_127-v0013} integral.}}{\placedmargin{Not automatically satisfied when the $X_s$ are \var{\deleted{\illegible{1}}}{}{deletion}{ega25_25_127-v0013} integral.}}{marginal note placed}{ega25_25_127-v0012}
\tnote{25-127-1}{The abbreviations alg, lin, syst and div mean algebraically, linear, system and divisors. The formula name Syst. lin. div is preserved exactly. In the last word the source alternatives are ``multiplicatinn'' and ``multiplication''; the English lexical reading remains marked as uncertain.}
